% =====================================================================
%  Copertina Lulu  --  comando \cover[chiave=valore]
%  Default: 6"x9" (152.4x228.6 mm), dorso 0.898" (372 pagine).
%  Le dimensioni sono in POLLICI (numeri puri).
% =====================================================================
\usepackage[T1]{fontenc}
\usepackage[utf8]{inputenc}
\usepackage{tikz}                 % carica anche pgfmath e pgfkeys
\usetikzlibrary{calc}
\pagestyle{empty}\setlength{\parindent}{0pt}

\usepackage{fontspec}

% ---------------------------------------------------------------------
%  Imposta la dimensione fisica della pagina (indipendente da geometry,
%  cosi' puo' essere decisa dal comando). Engine-agnostico.
% ---------------------------------------------------------------------
\newcommand{\cvrsetpage}[2]{%
    \setlength{\paperwidth}{#1}\setlength{\paperheight}{#2}%
    \ifdefined\pagewidth    \pagewidth=#1\relax    \pageheight=#2\relax    \fi
    \ifdefined\pdfpagewidth \pdfpagewidth=#1\relax \pdfpageheight=#2\relax \fi
}

% =====================================================================
%  Comando  \cover[ chiave=valore, ... ]
% =====================================================================
\newif\ifcvrguides   % linee guida on/off
\newif\ifcvrbc       % riquadro barcode on/off
\newsavebox{\cvrfcbox} % immagine del fronte
\pgfkeys{
    /cvr/.cd,
% --- geometria (pollici) ---
    trimw/.initial=6,          % larghezza trim
    trimh/.initial=9,          % altezza trim
    spinewidth/.initial=0.898, % larghezza dorso (NB: numero, non testo)
    bleed/.initial=0.125,      % abbondanza per lato
    safety/.initial=0.5,       % margine di sicurezza dal trim
    spinesafe/.initial=0.0625, % rientro sicurezza sul dorso
    bcw/.initial=3.622,        % larghezza area barcode
    bch/.initial=1.26,         % altezza area barcode
    bcup/.initial=0.5,         % distanza barcode dal taglio (trim) inferiore
% --- contenuti ---
    title/.initial=,           % titolo (fronte)
    subtitle/.initial=,        % sottotitolo (fronte)
    author/.initial=,          % autore (fronte, in basso)
    back/.initial=,            % testo di quarta (retro)
    imprint/.initial=,         % editore / sito (retro, sopra il barcode)
    spine/.initial=,           % testo del dorso (verticale)
    bgcolor/.initial=white,    % colore di fondo (bleed incluso)
    frontcover/.initial=,      % immagine sul fronte rifilato (senza bleed), sotto i testi
% --- interruttori ---
    barcode/.code={\csname cvrbc#1\endcsname},     % barcode=true|false
    guides/.code={\csname cvrguides#1\endcsname},  % guides=true|false
}
\newcommand{\cvrget}[1]{\pgfkeysvalueof{/cvr/#1}}

\newcommand{\cover}[1][]{%
% 1) default + chiavi utente (in due chiamate: evita il bug dell'a-capo)
    \cvrguidestrue\cvrbctrue
    \pgfkeys{/cvr/.cd, trimw=6, trimh=9, spinewidth=0.898, bleed=0.125,
        safety=0.5, spinesafe=0.0625, bcw=3.622, bch=1.26, bcup=0.5,
        title=, subtitle=, author=, back=, imprint=, spine=, bgcolor=white,
        frontcover=}%
    \pgfkeys{/cvr/.cd,#1}%
    % 2) coordinate derivate (origine in basso a sinistra)
    \pgfmathsetmacro\clDOCW{2*\cvrget{bleed}+2*\cvrget{trimw}+\cvrget{spinewidth}}
    \pgfmathsetmacro\clDOCH{2*\cvrget{bleed}+\cvrget{trimh}}
    \pgfmathsetmacro\clBACKL{\cvrget{bleed}}
    \pgfmathsetmacro\clBACKR{\cvrget{bleed}+\cvrget{trimw}}
    \pgfmathsetmacro\clSPL{\clBACKR}
    \pgfmathsetmacro\clSPR{\clBACKR+\cvrget{spinewidth}}
    \pgfmathsetmacro\clFRONTL{\clSPR}
    \pgfmathsetmacro\clFRONTR{\clSPR+\cvrget{trimw}}
    \pgfmathsetmacro\clTRIMB{\cvrget{bleed}}
    \pgfmathsetmacro\clTRIMT{\cvrget{bleed}+\cvrget{trimh}}
    \pgfmathsetmacro\clBACKSL{\clBACKL+\cvrget{safety}}
    \pgfmathsetmacro\clBACKSR{\clBACKR-\cvrget{safety}}
    \pgfmathsetmacro\clFRONTSL{\clFRONTL+\cvrget{safety}}
    \pgfmathsetmacro\clFRONTSR{\clFRONTR-\cvrget{safety}}
    \pgfmathsetmacro\clSAFEB{\clTRIMB+\cvrget{safety}}
    \pgfmathsetmacro\clSAFET{\clTRIMT-\cvrget{safety}}
    \pgfmathsetmacro\clSPSL{\clSPL+\cvrget{spinesafe}}
    \pgfmathsetmacro\clSPSR{\clSPR-\cvrget{spinesafe}}
    \pgfmathsetmacro\clBCR{\clBACKSR}
    \pgfmathsetmacro\clBCL{\clBACKSR-\cvrget{bcw}}
    \pgfmathsetmacro\clBCB{\clTRIMB+\cvrget{bcup}}
    \pgfmathsetmacro\clBCT{\clBCB+\cvrget{bch}}
    \pgfmathsetmacro\clTW{\cvrget{trimw}-2*\cvrget{safety}}  % larghezza utile
    \pgfmathsetmacro\clTH{\cvrget{trimh}-2*\cvrget{safety}}  % altezza utile
    % 3) dimensione fisica della pagina
    \cvrsetpage{\clDOCW in}{\clDOCH in}%
    % immagine del fronte: scalata per coprire il pannello rifilato (dal dorso al taglio
    % esterno, senza abbondanza) mantenendo le proporzioni; l'eccedenza viene tagliata
    \pgfkeysgetvalue{/cvr/frontcover}\cvrfc
    \ifx\cvrfc\empty\else
        \sbox{\cvrfcbox}{\includegraphics[width=\cvrget{trimw} in]{\cvrfc}}%
        \ifdim\dimexpr\ht\cvrfcbox+\dp\cvrfcbox\relax<\cvrget{trimh} in
            \sbox{\cvrfcbox}{\includegraphics[height=\cvrget{trimh} in]{\cvrfc}}%
        \fi
    \fi
    % 4) disegno
    \begin{tikzpicture}[x=1in, y=1in, remember picture, overlay]
    \coordinate (O) at (current page.south west);
    \begin{scope}[shift={(O)}]

    % sfondo: copre l'intero documento, bleed incluso
    \edef\cvrbg{\cvrget{bgcolor}}%
    \fill[fill=\cvrbg] (0,0) rectangle (\clDOCW,\clDOCH);

    % FRONTE: immagine sul pannello rifilato (l'abbondanza resta del colore di fondo)
    \ifx\cvrfc\empty\else
    \begin{scope}
    \clip (\clFRONTL,\clTRIMB) rectangle (\clFRONTR,\clTRIMT);
    \node[inner sep=0pt, anchor=center]
    at ($(\clFRONTL,\clTRIMB)!0.5!(\clFRONTR,\clTRIMT)$) {\usebox{\cvrfcbox}};
    \end{scope}
    \fi

    \ifcvrguides
    \draw[cyan!70,thick]   (0,0) rectangle (\clDOCW,\clDOCH);
    \draw[red,thick]       (\clBACKL,\clTRIMB) rectangle (\clFRONTR,\clTRIMT);
    \draw[magenta,thick]   (\clSPL,\clTRIMB) -- (\clSPL,\clTRIMT);
    \draw[magenta,thick]   (\clSPR,\clTRIMB) -- (\clSPR,\clTRIMT);
    \draw[blue!60,dashed]  (\clBACKSL,\clSAFEB)  rectangle (\clBACKSR,\clSAFET);
    \draw[blue!60,dashed]  (\clFRONTSL,\clSAFEB) rectangle (\clFRONTSR,\clSAFET);
    \draw[magenta!50,dashed](\clSPSL,\clSAFEB)   rectangle (\clSPSR,\clSAFET);
    \node[cyan!60!black,font=\scriptsize,anchor=north west]
    at (0.06,\clDOCH-0.06) {Documento con bleed};
    \node[red,font=\scriptsize,anchor=south]
    at ($(\clBACKL,\clTRIMT)!0.5!(\clSPL,\clTRIMT)$) {RETRO};
    \node[red,font=\scriptsize,anchor=south]
    at ($(\clSPR,\clTRIMT)!0.5!(\clFRONTR,\clTRIMT)$) {FRONTE};
    \fi

    \ifcvrbc
    \draw[orange,thick] (\clBCL,\clBCB) rectangle (\clBCR,\clBCT);
    \node[orange!80!black,font=\footnotesize]
    at ($(\clBCL,\clBCB)!0.5!(\clBCR,\clBCT)$) {AREA BARCODE};
    \fi

% FRONTE: titolo
    \node[anchor=north, text width=\clTW in, align=left]
    at ($(\clFRONTSL,\clSAFET)!0.5!(\clFRONTSR,\clSAFET)$)
        {\Huge \cvrget{title}\\[6pt]
    \normalfont\normalsize \cvrget{subtitle}};
    % FRONTE: autore
    \node[anchor=south, inner sep=0pt, text width=\clTW in, align=center]
    at ($(\clFRONTSL,\clSAFEB)!0.5!(\clFRONTSR,\clSAFEB)$)
        {\Large \cvrget{author}};
    % RETRO: quarta
    \node[anchor=north west, inner sep=0pt, text width=\clTW in, align=left]
    at (\clBACKSL,\clSAFET) {\cvrget{back}};
    % RETRO: editore
    \node[anchor=south west, inner sep=0pt, text width=1.6in, align=left]
    at (\clBACKSL,\clBCB) {\footnotesize \cvrget{imprint}};
    % DORSO
    \node[rotate=90, text width=\clTH in, align=flush center]
    at ($(\clSPL,\clTRIMB)!0.5!(\clSPR,\clTRIMT)$) {\cvrget{spine}};

    \end{scope}
    \end{tikzpicture}%
}

% fattore di scala rispetto al formato di riferimento 6x9
\newcommand{\cvrk}{\fpeval{\cvrget{trimw}/6}}
% \relativefontsize{corpo}{interlinea}: come \fontsize...\selectfont, ma scalato
\newcommand{\relativefontsize}[2]{\fontsize{\fpeval{#1*\cvrk}}{\fpeval{#2*\cvrk}}\selectfont}